% ID: MAT-TRI-EQG-010
% Titolo: Sistema goniometrico (2)
% Disciplina: Matematica
% Area: Trigonometria
% Argomento: Equazioni goniometriche
% Classe: Quarta secondaria di secondo grado
% Difficolta: 4
% Tipo: problema_numerico
% Tag: matematica, trigonometria, equazioni_goniometriche
% Fonte: verifica_goniometria_intermedia_anonimizzata
% Provenienza: materiale fornito dal docente; testo anonimizzato; soluzione redatta per Exergo

\begin{esercizio}
Risolvi il seguente sistema: 
\[
\begin{cases}
\sin 2x=-\sin x,\\[2mm]
\cos 2x=\cos x.
\end{cases}
\]
Salvo diversa indicazione, lavora in $\mathbb R$; angoli in radianti.
\end{esercizio}

\begin{soluzione}
La seconda equazione equivale a $(2\cos x+1)(\cos x-1)=0$.
Se $\cos x=1$ anche la prima è vera; se $\cos x=-1/2$, $\sin2x=-\sin x$.
Quindi $x=2k\pi$, $x=2\pi/3+2k\pi$ oppure $x=4\pi/3+2k\pi$. In tutte le famiglie di soluzioni $k\in\mathbb{Z}$.
\end{soluzione}
